% ============================================================================
% Fichier  : partie6/P06_C30_fondements_cryptanalyse.tex
% Titre    : Fondements de la Conception et Cryptanalyse
% Auteur   : MINKA MI NGUIDJOÏ Thierry Emmanuel
% Version  : 2.0 -- Juin 2026
% ============================================================================

\chapter{Fondements de la Conception et Cryptanalyse}
\label{chap:cryptanal}

\epigraph{%
    « La sécurité n'est pas un produit, mais un processus. »%
}{Bruce Schneier}

\niveauChapitre{Recherche}{%
    Cryptographie symétrique et asymétrique requises.
    Chapitres~\ref{chap:CRO} et~\ref{chap:prim-crypt} (Trilemme CRO
    formalisé et appliqué) recommandés~: ce chapitre relie la
    conception sécurisée et la cryptanalyse au cadre théorique déjà
    établi.%
}

\begin{boxobjectifs}
\begin{itemize}
    \item Maîtriser les quatre principes fondamentaux de la conception
          cryptographique sécurisée et leur justification.
    \item Relier le Trilemme CRO (Ch.~\ref{chap:CRO}) à la conception
          de protocoles~: comprendre en quoi il constitue une boussole
          de conception, pas seulement un outil d'analyse a posteriori.
    \item Classer rigoureusement les vulnérabilités cryptographiques
          selon une taxonomie à six catégories.
    \item Distinguer les modèles d'analyse \textit{black-box} et
          \textit{white-box} et savoir quand chacun s'applique.
    \item Comprendre les enjeux de la cryptanalyse post-quantique
          en lien avec le Chapitre~\ref{chap:quantique}.
    \item Adopter une posture critique et rigoureuse face aux
          publications scientifiques en cours, en distinguant
          résultat établi et prépublication non encore évaluée.
\end{itemize}
\end{boxobjectifs}

\bigskip

% ============================================================================
\section*{De l'analyse à la conception~: l'autre face du Trilemme}
% ============================================================================

Les Chapitres~\ref{chap:CRO} et~\ref{chap:prim-crypt} ont analysé
des protocoles \textbf{existants} à travers le cadre CRO~: comment
mesurer, a posteriori, le compromis confidentialité-fiabilité-
opposabilité d'Helios, d'un zk-SNARK, d'ECDSA. Ce chapitre inverse
la perspective~: comment \textbf{concevoir} un protocole en
anticipant ce compromis, et comment la \textbf{cryptanalyse},
la discipline qui teste la résistance d'un protocole à un adversaire
actif, constitue le complément indispensable de toute conception
sérieuse.

\separateur

% ============================================================================
\section{Philosophie de la Conception Sécurisée}
\label{sec:philosophie-conception}
% ============================================================================

\subsection{Quatre principes fondamentaux}

La conception de protocoles cryptographiques dignes de confiance
repose sur des principes qui anticipent délibérément la présence
d'un adversaire actif, et non simplement d'un environnement
défaillant.

\begin{boxdef}
\textbf{Devoir de méfiance} (\textit{Zero Trust})~: ne faire confiance
à aucun composant, message ou entité sans vérification préalable et
explicite. Aucune confiance implicite n'est accordée du seul fait
qu'un acteur se trouve ``à l'intérieur'' d'un périmètre supposé sûr.

\textbf{Minimalisme}~: réduire la surface d'attaque au strict
nécessaire. Toute fonctionnalité, tout code, toute complexité
supplémentaire constitue une opportunité additionnelle pour
l'attaquant, principe directement lié à la réduction de
$H(\Sigma)$ (Ch.~\ref{chap:CRO}), l'entropie de l'espace des preuves.

\textbf{Défense en profondeur}~: empiler plusieurs mécanismes de
sécurité indépendants, de sorte que la compromission d'une seule
couche n'entraîne pas l'effondrement de l'ensemble du système.

\textbf{\textit{Fail-safe defaults}}~: un système doit refuser
l'accès par défaut~; l'autorisation n'est accordée qu'explicitement,
après vérification réussie, jamais l'inverse.
\end{boxdef}

\begin{boxCRO}
\textbf{Analyse CRO. Les quatre principes relus à travers le
cadre formalisé}

Ces quatre principes, bien antérieurs à la formalisation du Trilemme
CRO, peuvent désormais être relus avec les outils des
Chapitres~\ref{chap:CRO} et~\ref{chap:prim-crypt}~:

Le \textbf{minimalisme} agit directement sur $H(\Sigma)$~: moins de
complexité signifie un espace de preuves plus restreint, donc
potentiellement plus facile à interpréter institutionnellement
(rappelons la borne~$\mathsf{Opp} \leq H(C)/H(V\mid C)$ du
Théorème~2, Ch.~\ref{chap:CRO}, où un espace de preuves trop riche
nuit à l'opposabilité si le contexte $H(C)$ ne suit pas).

La \textbf{défense en profondeur} est, formellement, une instance
de l'architecture en couches (Gold-Iron-Clay, Ch.~\ref{chap:CRO})~:
chaque couche apporte une force dominante distincte plutôt que de
chercher une primitive unique optimisant simultanément les trois
dimensions, ce que le théorème d'impossibilité interdit.

\cro{$\sim$ selon couche}{$\uparrow$ via minimalisme}{$\uparrow$ via espace de preuves restreint}
\end{boxCRO}

\subsection{Le Trilemme CRO comme boussole de conception}

Le Chapitre~\ref{chap:CRO} a établi, avec preuve mathématique
publiée, qu'\textbf{aucune construction ne peut maximiser
simultanément confidentialité, fiabilité et opposabilité} au-delà
d'un certain seuil conjoint (Théorème~3). Cette impossibilité n'est pas
une fatalité paralysante pour la conception, elle en est, au
contraire, la \textbf{boussole}~: une conception sécurisée ne
cherche pas l'optimum impossible, mais l'\textbf{équilibre
le mieux adapté au cas d'usage}, en suivant la méthode de sélection
de protocole établie au Ch.~\ref{chap:prim-crypt}
(Section~\ref{sec:implications-conception-prim})~: évaluer la
capacité institutionnelle cible, identifier le régime correspondant,
choisir l'architecture en conséquence.

\begin{boiteRemarque}{Sur les travaux de l'auteur en cours de
publication~: protocole ZK-NR}
Au-delà du cadre CRO (Ch.~\ref{chap:CRO}-\ref{chap:prim-crypt}),
l'auteur de ce cours développe, dans le cadre de travaux
de~recherche en cours, un protocole nommé \textit{ZK-NR}~\cite{minka2025zknr}
(\textit{Zero-Knowledge Non-Repudiation}), visant à combiner
plusieurs primitives (engagements de Merkle, preuves de type
STARK, signatures à seuil, signatures post-quantiques) pour des
scénarios de non-répudiation à forte criticité, une application
directe de la logique de défense en profondeur de ce chapitre.

\textbf{Statut à la date de rédaction de ce cours}~: contrairement
au cadre CRO (Ch.~\ref{chap:CRO}-\ref{chap:prim-crypt}), publié et
évalué par les pairs dans IEEE Access, le protocole ZK-NR ne fait
pas, à ce stade, l'objet d'un article complet disponible pour
vérification dans le cadre de ce cours. Conformément à la rigueur
méthodologique appliquée tout au long de cette Partie~VI, ce
chapitre s'abstient donc de présenter des détails techniques
(architecture précise, preuves de sécurité, scores de performance)
qui ne pourraient être vérifiés contre une publication accessible.
\textbf{Ce qui peut être affirmé en toute rigueur}~: le projet
illustre, par sa conception même (combinaison de primitives aux
forces complémentaires), le principe de défense en profondeur
appliqué à la non-répudiation, une instanciation conceptuelle
légitime à mentionner, indépendamment de la disponibilité d'un
article finalisé.
\end{boiteRemarque}

% ============================================================================
\section{Taxonomie des Failles Cryptographiques}
\label{sec:taxonomie-failles}
% ============================================================================

Comprendre l'attaquant nécessite de catégoriser systématiquement
ses vecteurs d'attaque possibles. Cette taxonomie, standard en
cryptographie, structure toute analyse de sécurité rigoureuse.

\begin{tableaucours}{p{3.5cm}p{7.5cm}}{%
    \tableauHeader{Type de faille} &
    \tableauHeader{Description et exemples}
}
\textbf{Conceptuelle (modèle)} & Faille dans la spécification
    formelle du protocole. L'attaquant respecte scrupuleusement
    le protocole tel que défini, mais en exploite une faiblesse
    logique. \textit{Exemple~:} attaque par rejeu de session,
    absence de protection contre la répétition
    (\textit{replay protection}). \\
\textbf{D'implémentation} & Faille dans le code source, malgré une
    spécification formelle correcte. \textit{Exemple~:} fuite de
    mémoire, gestion erronée des cas d'erreur, attaques par mesure
    de temps d'exécution (\textit{timing attacks}). \\
\textbf{Passive (écoute)} & L'adversaire observe uniquement, sans
    altérer la communication. \textit{Exemple~:} analyse de trafic
    (Ch.~\ref{chap:for-res}), cryptanalyse du texte chiffré
    intercepté. \\
\textbf{Active (modification)} & L'adversaire altère activement la
    communication. \textit{Exemple~:} attaque de l'homme du milieu
    (\textit{man-in-the-middle}), injection de messages forgés. \\
\textbf{Contre les primitives} & Attaque visant directement le
    fondement mathématique de la primitive. \textit{Exemple~:}
    algorithme de Shor contre RSA/ECC (Ch.~\ref{chap:quantique}),
    attaque par canal auxiliaire (consommation électrique,
    rayonnement électromagnétique) sur une implémentation ECC. \\
\textbf{Contre le protocole} & Attaque exploitant l'interaction
    entre plusieurs primitives correctement implémentées
    individuellement. \textit{Exemple~:} attaque par entrelacement
    (\textit{interleaving}) entre deux sessions, confusion des
    rôles entre prouveur et vérificateur. \\
\end{tableaucours}
\vspace{-0.8em}
\begin{center}{\small\itshape Tableau~30.1. Taxonomie des failles
de sécurité cryptographique}\end{center}

\begin{boxscenario}{Application de la taxonomie~: l'affaire MVOM
revisitée}
Reprenons l'analyse de l'Opération MVOM (Ch.~\ref{chap:for-sys},
\ref{chap:for-res}) à travers cette taxonomie~:

\begin{itemize}
    \item L'\textbf{accès SSH par force brute} relève d'une faille
          \textbf{active}~: l'attaquant ne se contente pas
          d'observer, il tente activement des authentifications.
    \item L'utilisation du \textbf{chiffrement BitLocker} sur le
          disque saisi, si elle avait été contournée par une faille
          dans l'implémentation Windows plutôt que par obtention
          légale de la clé, aurait relevé d'une faille
          \textbf{d'implémentation}, distincte d'une attaque
          \textbf{contre la primitive} AES elle-même
          (Ch.~\ref{chap:quantique}~: AES reste robuste face aux
          capacités actuelles).
    \item La \textbf{tentative \texttt{StopLogging}}
          (Ch.~\ref{chap:anti-for}) n'exploite aucune faille
          cryptographique~: elle relève d'une catégorie distincte,
          hors taxonomie cryptographique stricte, mais analogue
          conceptuellement à une faille \textbf{conceptuelle}~: le
          système de logging n'anticipait pas, dans sa conception
          initiale, qu'un attaquant disposant de privilèges
          administrateur chercherait à désactiver l'audit
          lui-même.
\end{itemize}
\end{boxscenario}

% ============================================================================
\section{Introduction à la Cryptanalyse}
\label{sec:intro-cryptanalyse}
% ============================================================================

\subsection{Objet et finalité}

\begin{boxdef}
\textbf{Cryptanalyse}~: discipline scientifique consistant à
étudier et tenter de contourner les mécanismes de protection
cryptographique d'un système. Son objectif \textbf{n'est pas
exclusivement offensif}~: elle est indispensable à la validation
de la solidité des constructions, exactement comme la falsification
active (E9bis du PIU, Ch.~6~(UE~SEC4031)) est indispensable à la
validation d'une hypothèse investigative, une hypothèse, ou un
protocole, qui n'a jamais été sérieusement attaqué n'a pas démontré
sa robustesse~; il n'a simplement pas encore été testé.
\end{boxdef}

Cette analogie entre cryptanalyse et falsification active n'est pas
fortuite~: les deux disciplines partagent une même épistémologie
poppérienne, la confiance se construit par résistance répétée
à des tentatives sérieuses de réfutation, pas par l'absence de
telles tentatives.

\subsection{Approches \textit{black-box} et \textit{white-box}}

\begin{boxdef}
\textbf{Analyse \textit{black-box}}~: l'attaquant ne dispose que
des entrées et sorties observables du système. Il déduit les
vulnérabilités potentielles par observation du comportement
externe, temps de réponse, consommation énergétique, messages
d'erreur.

\textbf{Analyse \textit{white-box}}~: l'attaquant dispose d'un
accès total à l'implémentation, au code source, voire aux données
internes en cours d'exécution. C'est le scénario le plus
défavorable pour le défenseur, et le standard d'évaluation pour
tout système fortement exposé.
\end{boxdef}

\begin{tableaucours}{p{2.8cm}p{4.0cm}p{4.5cm}}{%
    \tableauHeader{Modèle} &
    \tableauHeader{Information disponible} &
    \tableauHeader{Pertinence forensique (ce cours)}
}
\textit{Black-box} & Entrées/sorties observables uniquement &
    Analyse réseau (Ch.~\ref{chap:for-res})~: l'investigateur,
    face à un trafic chiffré inconnu, raisonne en \textit{black-box}, métadonnées, motifs de connexion, sans accès au contenu \\
\textit{White-box} & Code source, mémoire, état interne complet &
    Analyse de firmware (Ch.~\ref{chap:for-iot}, Binwalk)~:
    l'investigateur dispose d'un accès total au binaire extrait,
    scénario \textit{white-box} par excellence \\
\end{tableaucours}
\vspace{-0.8em}
\begin{center}{\small\itshape Tableau~30.2. Les deux modèles
appliqués à la pratique forensique de ce cours}\end{center}

\begin{boiteRemarque}{Le principe de Kerckhoffs, fondement du
modèle white-box}
L'adoption du modèle \textit{white-box} comme standard d'évaluation
en cryptographie moderne découle directement du \textbf{principe de
Kerckhoffs} (1883)~: la sécurité d'un système cryptographique doit
reposer uniquement sur le secret de la clé, jamais sur le secret
de l'algorithme. Un protocole dont la sécurité dépendrait de la
confidentialité de sa conception échouerait, par construction,
face à un adversaire \textit{white-box}, exactement le cas
d'un firmware d'appareil IoT analysé après extraction
(Ch.~\ref{chap:for-iot}), où l'algorithme complet est accessible
à l'investigateur (et potentiellement à l'attaquant qui aurait
pu l'extraire de la même façon).
\end{boiteRemarque}

% ============================================================================
\section{Cryptanalyse Post-Quantique~: Anticiper Deux Horizons Temporels}
\label{sec:cryptanalyse-pq}
% ============================================================================

Le Chapitre~\ref{chap:quantique} a détaillé les algorithmes de Shor
et Grover et leurs implications pour les preuves numériques. Cette
section les replace dans le cadre plus large de la
\textbf{méthodologie cryptanalytique}~: une analyse de sécurité
moderne ne peut plus se limiter à un seul horizon temporel.

\begin{enumerate}
    \item \textbf{Aujourd'hui}~: résistance aux attaques classiques
          sur les capacités de calcul actuellement disponibles, l'horizon d'analyse traditionnel de la cryptanalyse.
    \item \textbf{Demain}~: résistance aux attaques quantiques,
          notamment via les algorithmes de \textbf{Shor} (rupture
          de la cryptographie asymétrique) et de \textbf{Grover}
          (accélération quadratique réduisant la sécurité effective
          des clés symétriques), horizon détaillé au
          Ch.~\ref{chap:quantique}, Sections~1 et~3.
\end{enumerate}

\begin{boxCRO}
\textbf{Analyse CRO. ``Harvest Now, Decrypt Later'' comme
problème cryptanalytique prospectif}

La stratégie \textit{Harvest Now, Decrypt Later}
(Ch.~\ref{chap:quantique}, section~26.2) transforme la cryptanalyse
d'une discipline rétrospective (``ce protocole est-il cassé~?'')
en discipline \textbf{prospective} (``ce protocole sera-t-il
cassé, et quand~?''), rendant la cryptanalyse prospective
absolument critique pour la protection des secrets à très longue
durée de sensibilité, exactement comme établi au
Ch.~\ref{chap:quantique} pour les données de santé ou les
secrets d'État.

$\mathcal{C}$~: cette stratégie inverse la temporalité habituelle
de l'analyse de risque, la confidentialité d'aujourd'hui dépend
d'une cryptanalyse qui n'existe pas encore.

$\mathcal{R}$~: aucune méthode de cryptanalyse actuelle ne peut
\textbf{prouver} la résistance future d'un protocole~; elle ne peut
que documenter l'absence de faille connue \emph{à ce jour}, une
limite épistémologique fondamentale, déjà signalée au
Ch.~\ref{chap:quantique} concernant les calendriers spéculatifs.

\cro{$\downarrow\downarrow$ pour données longue durée}{$\sim$ limite épistémologique assumée}{$\sim$}
\end{boxCRO}

% ============================================================================
\section{De la Cryptanalyse à la Conception~: une Boucle Vertueuse}
\label{sec:boucle-vertueuse}
% ============================================================================

\subsection{Le cycle conception-cryptanalyse-révision}

La pratique mature de la cryptographie institutionnelle, telle
qu'elle ressort des Chapitres~\ref{chap:CRO} et~\ref{chap:prim-crypt},
suit un cycle itératif~:

\begin{enumerate}
    \item \textbf{Concevoir} un protocole en appliquant les quatre
          principes de la Section~\ref{sec:philosophie-conception}
          et en explicitant, dès la conception, le régime de
          capacité institutionnelle visé (Ch.~\ref{chap:prim-crypt}).
    \item \textbf{Soumettre} le protocole à une cryptanalyse
          rigoureuse, \textit{white-box} de préférence
          (Section~\ref{sec:intro-cryptanalyse}), analogue à la
          revue par les pairs qui a validé les Théorèmes~1 à~3 du
          Ch.~\ref{chap:CRO} avant publication.
    \item \textbf{Documenter} les failles découvertes selon la
          taxonomie de la Section~\ref{sec:taxonomie-failles}, en
          distinguant leur catégorie précise.
    \item \textbf{Réviser} la conception en réponse, en
          réévaluant le positionnement CRO du protocole révisé
          (méthode du Ch.~\ref{chap:prim-crypt}).
    \item \textbf{Republier} ou redéployer, avec une documentation
          transparente du processus, exactement la posture de
          transparence méthodologique pratiquée par les auteurs
          du cadre CRO eux-mêmes, qui identifient explicitement
          les limites et pistes futures de leur propre travail.
\end{enumerate}

\subsection{Application au contexte forensique}

\begin{boxjuris}
\textbf{Cryptanalyse et admissibilité en droit camerounais}

Un expert judiciaire camerounais (Ch.~12~(UE~SEC4031), experts
agréés Décret~69/DF/544) appelé à témoigner sur la solidité d'un
mécanisme cryptographique utilisé dans un dossier, par exemple,
la robustesse d'une signature électronique contestée
(Ch.~11~(UE~SEC4031), eIDAS), doit être en mesure de
répondre, en contre-interrogatoire (Ch.~\ref{chap:rapport}), à des
questions de cryptanalyse~: \textit{``Cet algorithme a-t-il fait
l'objet d'une cryptanalyse publiée~? A-t-il résisté aux meilleures
tentatives de la communauté scientifique~? Existe-t-il des attaques
connues, même théoriques, contre cette construction~?''} La
réponse attendue, conforme à la rigueur de ce chapitre, distingue
toujours~: ce qui est \textbf{publié et vérifié} (par exemple, les
théorèmes du Ch.~\ref{chap:CRO}), ce qui relève d'une
\textbf{recherche en cours} non encore validée par les pairs (par
exemple, le protocole ZK-NR évoqué en
Section~\ref{sec:philosophie-conception}), et ce qui demeure une
\textbf{question ouverte}.
\end{boxjuris}

\separateur

\begin{boxresume}
\textbf{Ce qu'il faut retenir de ce chapitre~:}
\begin{itemize}
    \item \textbf{Quatre principes de conception}~: Zero Trust,
          minimalisme, défense en profondeur, \textit{fail-safe
          defaults}, chacun directement reliable au cadre CRO
          formalisé (Ch.~\ref{chap:CRO}).

    \item Le \textbf{Trilemme CRO} fonctionne comme \textbf{boussole
          de conception}, pas seulement comme outil d'analyse a
          posteriori~: il guide le choix d'architecture vers
          l'équilibre adapté au régime de capacité institutionnelle
          cible (Ch.~\ref{chap:prim-crypt}).

    \item \textbf{Taxonomie en six catégories}~: failles
          conceptuelles, d'implémentation, passives, actives,
          contre les primitives, contre le protocole, grille
          systématique applicable à toute analyse de sécurité,
          y compris en contexte forensique.

    \item \textbf{Black-box vs white-box}~: deux modèles d'analyse
          directement transposables à la pratique forensique de ce
          cours (analyse réseau~$\approx$ black-box~; analyse de
          firmware~$\approx$ white-box), enracinés dans le principe
          de Kerckhoffs.

    \item \textbf{Cryptanalyse post-quantique}~: deux horizons
          temporels à anticiper simultanément (aujourd'hui/demain),
          avec la stratégie \textit{Harvest Now, Decrypt Later}
          transformant la cryptanalyse en discipline prospective
          (Ch.~\ref{chap:quantique}).

    \item \textbf{Cycle vertueux conception-cryptanalyse-révision}~:
          la pratique mature ne sépare pas concevoir et attaquer~;
          elle les itère, avec une documentation transparente
          à chaque étape.

    \item \textbf{Rigueur sur les travaux en cours}~: la distinction
          entre résultat publié et évalué par les pairs (Trilemme
          CRO), prépublication non encore évaluée par les pairs (ZK-NR), et
          question ouverte, est elle-même une compétence
          cryptanalytique et forensique, appliquée explicitement
          dans ce chapitre à propos des propres travaux de l'auteur.
\end{itemize}
\end{boxresume}

\bigskip

\begin{boxexo}[title={\ding{45}\quad Exercice~30.1. Application de la taxonomie \niveaubase}]
Pour chacune des situations suivantes, rencontrées dans ce cours,
classez la faille selon la taxonomie du Tableau~30.1 (une ou
plusieurs catégories peuvent s'appliquer, justifiez)~:
\begin{enumerate}[(a)]
    \item Un attaquant intercepte du trafic réseau chiffré sans
          le modifier, dans l'espoir de le déchiffrer plus tard
          (Ch.~\ref{chap:quantique}, \textit{Harvest Now, Decrypt
          Later}).
    \item Un développeur laisse une clé API en clair dans une
          couche Docker, supprimée dans une couche ultérieure mais
          toujours accessible dans l'historique
          (Ch.~\ref{chap:containers}).
    \item Un attaquant exploite l'algorithme de Shor contre une
          signature RSA (Ch.~\ref{chap:quantique}).
    \item Un acteur malveillant intercale ses propres messages
          dans un échange SSH déjà établi.
\end{enumerate}
\end{boxexo}

\begin{boxexo}[title={\ding{45}\quad Exercice~30.2. Conception sous contrainte CRO \niveaumoyen}]
Vous devez concevoir un mécanisme de signature pour les rapports
d'expertise forensique numérique au Cameroun (Ch.~\ref{chap:rapport}),
destiné à être présenté devant des tribunaux de juridiction
générale ($H(C) \approx 8$~bits selon le Ch.~\ref{chap:prim-crypt}).
\begin{enumerate}[(a)]
    \item En appliquant la méthode de sélection de protocole
          (Ch.~\ref{chap:prim-crypt}, section~28.8.1), dans quel
          régime de capacité institutionnelle ce contexte se
          situe-t-il~? Quel type de protocole (transparent ou
          préservant la confidentialité) recommanderiez-vous~?
    \item Appliquez les quatre principes de conception
          (Section~\ref{sec:philosophie-conception}) à ce
          mécanisme~: que signifierait concrètement le
          ``minimalisme'' pour une signature de rapport d'expertise~?
    \item Quelle catégorie de faille (Tableau~30.1) vous semble la
          plus critique à anticiper pour ce cas d'usage spécifique~?
          Justifiez et proposez une contre-mesure de conception.
\end{enumerate}
\end{boxexo}

\begin{boxexo}[title={\ding{45}\quad Exercice~30.3. Esprit critique face à un protocole en développement \niveauexpert}]
Un consortium international vous présente un nouveau protocole de
non-répudiation pour la chaîne de custody numérique, accompagné
d'une présentation affirmant~: ``Notre protocole atteint
simultanément confidentialité~$=0{,}95$, fiabilité~$=0{,}98$ et
opposabilité~$=0{,}92$, soit un score combiné de~$0{,}857$, le
plus élevé jamais documenté.''
\begin{enumerate}[(a)]
    \item En vous appuyant sur le Théorème~3 (Trilemme CRO
          conditionnel, Ch.~\ref{chap:CRO}) et sur les valeurs
          réellement publiées du Tableau~28.3 (Ch.~\ref{chap:prim-crypt},
          produits CLO entre~$0{,}27$ et~$0{,}57$), évaluez la
          plausibilité de cette affirmation.
    \item Quelles questions précises poseriez-vous à ce consortium
          pour vérifier la validité de ces chiffres (référez-vous
          à la méthodologie en quatre étapes du
          Ch.~\ref{chap:prim-crypt}, section~28.1)~?
    \item Ce chapitre a traité le protocole ZK-NR, mentionné dans
          le Prologue de ce cours, avec une prudence méthodologique
          explicite face à l'absence d'article complet vérifiable.
          Rédigez un court paragraphe expliquant pourquoi cette
          même prudence devrait s'appliquer à l'affirmation du
          consortium de cet exercice, en distinguant ce qui
          relèverait d'un résultat publié et évalué par les pairs
          de ce qui relève d'une affirmation commerciale non
          vérifiée.
\end{enumerate}
\end{boxexo}

\bigskip

\begin{boxinfo}
\textbf{Transition.} Ce chapitre clôt l'exploration de la
conception et de la cryptanalyse fondées sur le cadre CRO. Le
chapitre suivant aborde la \textbf{méthodologie d'analyse formelle}~:
comment structurer rigoureusement une preuve de sécurité, et
comment cette même rigueur méthodologique s'applique à la rédaction
des preuves forensiques présentées tout au long de ce cours.
\end{boxinfo}